\documentclass[runningheads]{llncs}

\usepackage{amsmath,amssymb}
\usepackage{graphicx}
\usepackage{booktabs}
\usepackage{tabularx}
\usepackage{array}
\usepackage{url}
\usepackage{hyperref}
\usepackage{enumitem}

\newcommand{\pending}[1]{\par\noindent\textbf{[AWAITING HUMAN RESULTS: #1]}\par}

\begin{document}

\title{MESA: A Meta-Evaluation Standard and Audit Framework for AI Benchmark Requirements}
\titlerunning{MESA: AI Benchmark Audit for Requirements Engineering}

\author{Anonymous Author(s)}
\authorrunning{Anonymous Author(s)}

\institute{Affiliation Withheld for Double-Blind Review\\
\email{contact@example.org}}

\maketitle

\begin{abstract}
\textbf{Context and motivation:} In Requirements Engineering for Artificial Intelligence (RE4AI), benchmark scores are frequently cited as acceptance evidence to verify that an AI component satisfies system capability requirements. However, using benchmark scores as sound requirements evidence requires verifying that benchmark tasks, scoring implementations, and documentation actually support the claimed capabilities.

\textbf{Question/problem:} Existing benchmark review frameworks lack common criteria, use divergent scales, and conflate missing documentation with system defects, making it difficult for software engineering teams to evaluate benchmark validity or pool reviews. Furthermore, LLM-based audit tools can miss code defects or hallucinate checks if not audited by human engineers.

\textbf{Principal ideas/results:} We introduce MESA (Meta-Evaluation Standard for AI), a comparative standard mapping 280 items from seven review frameworks onto 30 shared criteria, decoupling unstated practices from verified defects. We evaluate the derived MESA rubric through 28 coding agent sessions and a human-in-the-loop audit platform where software engineers verify agent-generated ratings and code locators for HumanEval. Across 143 cross-framework disagreements, 92 stem from scoring rules rather than evidence differences. In human verification, pairing engineers with agent evidence records cut review time to $\sim$32 minutes while enabling auditors to detect live code execution defects that paper-only reading missed.

\textbf{Contribution:} MESA provides a requirements metrology framework and human-agent audit methodology for verifying AI benchmark evidence, ensuring valid acceptance criteria for AI-enabled software systems.

\keywords{Requirements Engineering for AI \and AI Benchmarks \and Acceptance Testing \and Construct Validity \and Human-in-the-Loop Audit \and Empirical Software Engineering.}
\end{abstract}

\section{Introduction}
\label{sec:introduction}

Software engineering teams specifying requirements for artificial intelligence (AI) systems face an evidential gap: when selecting a commercial or open-weight AI model to satisfy functional requirements (such as code generation correctness or factual question answering), teams routinely rely on published benchmark scores as verification evidence \cite{habiba_re4ai_maturity_2024}. A high score on a public leaderboard, however, does not establish that an AI component satisfies deployment requirements. Benchmark scores depend heavily on task construction, execution environments, scoring heuristics, and baseline comparisons. Leaderboard positions often encourage generalizations that exceed the underlying evidence \cite{singh_leaderboard_illusion,alonso_church_sigir}. In requirements engineering (RE), rigorous verification demands inspecting the benchmark itself: do its tasks, test harness, and scoring implementation validly operationalize the required capability?

Since 2023, at least seven frameworks have been proposed to review AI benchmarks as measurement instruments \cite{reuel_betterbench,liu_ecbd_2024,san_joaquin_scorecard_2026,cao_how2bench_2025,zhu_agentic_benchmark_checklist_2025,corona_benchmark_caliper_2026,cohn_hernandez_orallo_oecd_2023}. When requirements engineers attempt to apply them, they encounter an immediate obstacle: existing frameworks lack common criteria, use incompatible scoring scales, and enforce conflicting rules for unstated information. For example, when a benchmark paper omits contamination safeguards, BetterBench scores 0, the Scorecard marks the requirement unsatisfied, and ECBD records the absence in text without a penalty. An engineer comparing reviews cannot distinguish whether a low score reflects a confirmed defect or merely divergent scoring conventions. Furthermore, negative scores for unstated details outnumber documented defects by over ten to one, penalizing incomplete reporting rather than empirical failure.

To bridge this gap, we introduce \textbf{MESA (Meta-Evaluation Standard for AI)}, a translation layer and audit standard for benchmark reviews. MESA establishes:
\begin{enumerate}[leftmargin=*]
    \item \textbf{30 Shared Criteria} mapping all 280 items from the seven published frameworks onto a common metrology;
    \item \textbf{Three Standardized Verdicts} (\emph{satisfied}, \emph{unsatisfied}, \emph{open}) with qualifiers isolating missing documentation (\emph{undocumented}) from confirmed defects;
    \item \textbf{The Derived MESA Rubric}, comprising 8 descriptive facets (Part~A) and 29 evaluative questions (Part~B) with a six-anchor evidence scale;
    \item \textbf{A Human-in-the-Loop Audit Workbench}, enabling software engineering teams to audit LLM-generated benchmark reviews, trace claims to repository code, and catch reasoning flaws.
\end{enumerate}

We evaluate MESA through three empirical studies: (1) a cross-framework comparison applying six frameworks to four benchmarks (HumanEval, TruthfulQA, Humanity's Last Exam, and $\tau$-bench) using coding agents; (2) a repeatability study of the derived MESA rubric across two runs; and (3) an empirical human-in-the-loop audit study with software engineers auditing HumanEval using our interactive workbench.

\section{Background and Related Work}
\label{sec:background}

\subsection{AI Benchmark Scores as Requirements Evidence}
In Requirements Engineering for AI (RE4AI), establishing acceptance criteria for probabilistic, data-dependent components remains a primary challenge \cite{habiba_re4ai_maturity_2024}. When software requirements state that a model must perform docstring-to-code synthesis (e.g. Copilot integration), engineers require traceability between the specified requirement, the benchmark tasks, the execution guardrails, and the acceptance decision \cite{bean_measuring_what_matters}. Weak construct representation, data contamination, and fragile prompts compromise this evidential chain \cite{jo_wilson_what_does_your_benchmark_measure,line_goes_up}. In traditional software testing, test cases are audited for boundary coverage and oracle correctness; similarly, AI benchmarks serving as requirement oracles must undergo systematic audit before acceptance.

\subsection{Construct Validity and Test Standards}
In educational and psychological measurement, validity is defined as the degree to which theory and empirical evidence support the interpretation of test scores for a specified use \cite{kane_validation,aera_apa_ncme_standards_2014}. The European Federation of Psychologists' Associations (EFPA) Test Review Model 2025 operationalizes this by separating neutral instrument description (Part~A) from evidence evaluation (Part~B) \cite{efpa_test_review_model_2025}. Cohn and Hern\'andez-Orallo adapted this taxonomic approach to AI evaluation across 18 design facets \cite{cohn_hernandez_orallo_oecd_2023}. MESA extends this metrology to benchmark review for requirements engineering, ensuring that benchmark audits evaluate whether scores justify specific functional and safety claims.

\subsection{Existing Benchmark Review Frameworks}
We examine seven review frameworks: \textbf{BetterBench (BB)} \cite{reuel_betterbench} (47 lifecycle items), \textbf{ECBD} \cite{liu_ecbd_2024} (20 design questions), the \textbf{Scorecard (SC)} \cite{san_joaquin_scorecard_2026} (68 validity items), \textbf{How2Bench (H2B)} \cite{cao_how2bench_2025} (55 checklist items for code benchmarks), the \textbf{Agentic Benchmark Checklist (ABC)} \cite{zhu_agentic_benchmark_checklist_2025} (44 items), \textbf{Benchmark Caliper (CAL)} \cite{corona_benchmark_caliper_2026} (28 transfer items), and the \textbf{OECD facets} \cite{cohn_hernandez_orallo_oecd_2023} (18 facets). As shown in Table~\ref{tab:frameworks}, each framework addresses overlapping concerns with incompatible scoring conventions.

\begin{table}[t]
\caption{The seven review frameworks and the MESA rubric.}
\label{tab:frameworks}
\centering
\scriptsize
\setlength{\tabcolsep}{4pt}
\begin{tabularx}{\textwidth}{@{}l>{\raggedright\arraybackslash}p{2.2cm}rr>{\raggedright\arraybackslash}X>{\raggedright\arraybackslash}p{2.8cm}@{}}
\toprule
\textbf{Code} & \textbf{Framework} & \textbf{Items} & \textbf{Scope} & \textbf{Rule for Undocumented Items} & \textbf{Standardized Verdict} \\
\midrule
BB & BetterBench \cite{reuel_betterbench} & 47 & General AI & Score 0 (unacknowledged); Score 5 (no proof) & Unsatisfied, \emph{undocumented} \\
ECBD & ECBD \cite{liu_ecbd_2024} & 20 & NLP & Recorded in text without score & Open \\
SC & Scorecard \cite{san_joaquin_scorecard_2026} & 68 & General AI & Vague or absent documentation fails item & Unsatisfied, \emph{undocumented} \\
H2B & How2Bench \cite{cao_how2bench_2025} & 55 & Code & Absent documentation grouped with not done & Unsatisfied, \emph{undocumented} \\
ABC & ABC \cite{zhu_agentic_benchmark_checklist_2025} & 44 & Agentic & Undocumented item scored 0 & Unsatisfied, \emph{undocumented} \\
CAL & Caliper \cite{corona_benchmark_caliper_2026} & 28 & Cultural & Unknown; dimension lowered to 1--2 & Unsatisfied, \emph{undocumented} \\
OECD & OECD \cite{cohn_hernandez_orallo_oecd_2023} & 18 & All AI & N/A allowed on Reliability; others require value & Reliability open; others evaluated \\
\midrule
--- & \textbf{MESA Rubric} & 37 & General AI & Recorded as \emph{Undocumented} (open question) & Open \\
\bottomrule
\end{tabularx}
\end{table}

\section{The MESA Meta-Evaluation Standard and Rubric}
\label{sec:mesa_standard}

\subsection{Shared Review Criteria}
All 280 items across the seven frameworks were transcribed and mapped bottom-up onto \textbf{30 shared criteria} (C01 to C30). Table~\ref{tab:coverage} lists the 30 criteria and the number of items allocated by each framework. Framework coverage varies dramatically: Caliper concentrates all 28 items on cultural transfer (C23), whereas BetterBench, Scorecard, and How2Bench distribute items broadly across lifecycle and validity concerns.

\begin{table}[t]
\caption{The 30 shared review criteria and framework item allocations (RQ1).}
\label{tab:coverage}
\centering
\scriptsize
\setlength{\tabcolsep}{2.5pt}
\begin{tabular}{@{}llrrrrrrrr@{}}
\toprule
\textbf{Id} & \textbf{Shared Review Criterion} & \textbf{BB} & \textbf{ECBD} & \textbf{SC} & \textbf{H2B} & \textbf{ABC} & \textbf{CAL} & \textbf{OECD} & \textbf{Total} \\
\midrule
C01 & Ability under test is stated and concrete & 2 & 2 & 3 & 2 & 1 & 0 & 1 & 6 \\
C02 & Intended use, users and real-world relevance & 2 & 3 & 1 & 1 & 1 & 0 & 2 & 6 \\
C03 & Task selection justified for claimed ability & 4 & 3 & 5 & 2 & 1 & 0 & 0 & 5 \\
C04 & Task set covers the ability representatively & 2 & 5 & 4 & 5 & 0 & 5 & 1 & 6 \\
C05 & Tasks and reference answers were checked & 1 & 0 & 10 & 7 & 12 & 3 & 0 & 5 \\
C06 & Scoring rule is described & 3 & 2 & 4 & 4 & 3 & 0 & 1 & 6 \\
C07 & Scoring implementation checked / cannot be gamed & 1 & 1 & 8 & 1 & 18 & 0 & 0 & 5 \\
C08 & Aggregation and subscores are justified & 0 & 3 & 3 & 1 & 0 & 1 & 1 & 5 \\
C09 & Comparison points and baselines justified & 3 & 0 & 3 & 0 & 2 & 0 & 1 & 4 \\
C10 & Variation on repeated testing reported & 0 & 0 & 2 & 3 & 1 & 0 & 2 & 4 \\
C11 & Prompting and setup described and tested & 2 & 3 & 3 & 6 & 0 & 0 & 0 & 4 \\
C12 & Uncertainty and statistical support reported & 1 & 2 & 3 & 1 & 2 & 0 & 0 & 5 \\
C13 & Shortcuts, confounds and chance investigated & 0 & 1 & 12 & 1 & 6 & 2 & 1 & 6 \\
C14 & Training data contamination assessed & 2 & 0 & 1 & 1 & 2 & 0 & 0 & 4 \\
C15 & Rules for test training / repeat submissions & 1 & 0 & 0 & 0 & 1 & 0 & 0 & 2 \\
C16 & Claims beyond test supported (generalization) & 1 & 1 & 6 & 1 & 0 & 1 & 2 & 6 \\
C17 & Fairness across systems and conditions & 0 & 1 & 1 & 0 & 0 & 1 & 2 & 4 \\
C18 & Materials and replication instructions provided & 10 & 0 & 3 & 11 & 3 & 0 & 1 & 5 \\
C19 & Terms of access and license stated & 1 & 0 & 0 & 2 & 0 & 0 & 0 & 2 \\
C20 & Versions and changes are recorded & 3 & 0 & 1 & 0 & 3 & 0 & 2 & 4 \\
C21 & Limitations and flaws stated; claims bounded & 4 & 0 & 6 & 0 & 4 & 0 & 0 & 3 \\
C22 & Maintenance and contacts exist & 4 & 0 & 0 & 0 & 0 & 0 & 0 & 1 \\
C23 & Cultural and regional transfer assessed & 0 & 0 & 0 & 0 & 0 & 28 & 1 & 2 \\
C24 & Harmful or sensitive content handled & 1 & 0 & 0 & 3 & 0 & 1 & 0 & 3 \\
C25 & Evaluated systems and difficulty examined & 0 & 0 & 0 & 6 & 0 & 0 & 1 & 2 \\
C26 & Results presented and explained & 0 & 0 & 0 & 6 & 0 & 0 & 0 & 1 \\
C27 & Data provenance and collection documented & 7 & 1 & 4 & 5 & 0 & 1 & 2 & 6 \\
C28 & Execution environment isolated and controlled & 0 & 0 & 3 & 0 & 7 & 0 & 1 & 3 \\
C29 & Domain experts / users involved in review & 1 & 1 & 2 & 0 & 0 & 2 & 0 & 4 \\
C30 & Benchmark instrument descriptive classification & 0 & 0 & 0 & 0 & 0 & 0 & 10 & 1 \\
\bottomrule
\end{tabular}
\end{table}

\subsection{Standardized Verdicts and Qualifiers}
To compare reviews fairly, answers are standardized into three verdicts:
\begin{itemize}[leftmargin=*]
    \item \textbf{Satisfied (\emph{favorable}):} The benchmark meets the criterion standard.
    \item \textbf{Unsatisfied (\emph{unfavorable}):} The benchmark exhibits a confirmed defect or is scored failing.
    \item \textbf{Open (\emph{unresolved}):} Missing documentation, uncertainty, or inapplicability.
\end{itemize}
Crucially, when a framework penalizes silence by design, MESA attaches the \textbf{\emph{undocumented}} qualifier: the verdict reflects the framework's scoring rule, but identifies that the negative outcome stems from unstated information rather than a confirmed flaw.

\subsection{The Derived MESA Rubric}
Derived directly from the 30 criteria, the MESA rubric has two parts:
\begin{itemize}[leftmargin=*]
    \item \textbf{Part~A (8 OECD Facets):} Purpose, Judgeability, Reference, Containedness, Objectivity, Autonomy, Progression, and Variation. Three facets formally gate Part~B (e.g. Judgeability gates Question B07).
    \item \textbf{Part~B (29 Evaluative Questions):} Each question has an asymmetric six-anchor evidence scale: \emph{Claimed} (asserted without inspectable artifact), \emph{Documented} (inspectable artifact present), \emph{Verified} (systematic check reported with coverage/limits), \emph{Contradicted} (confirmed defect or contradiction), \emph{Undocumented} (insufficient evidence), and \emph{Not applicable}. Table~\ref{tab:rubric_summary} defines the anchor thresholds for key criteria.
\end{itemize}

\begin{table}[t]
\caption{Representative MESA Rubric Questions and Anchor Thresholds.}
\label{tab:rubric_summary}
\centering
\scriptsize
\setlength{\tabcolsep}{3pt}
\begin{tabularx}{\textwidth}{@{}l>{\raggedright\arraybackslash}p{2.6cm}>{\raggedright\arraybackslash}X>{\raggedright\arraybackslash}X>{\raggedright\arraybackslash}p{2.2cm}@{}}
\toprule
\textbf{Id} & \textbf{Evaluation Question} & \textbf{What Counts as Shown (Documented)} & \textbf{What Counts as Checked (Verified)} & \textbf{Applies When} \\
\midrule
B01 & Ability under test stated & Operational form and scope defined & Definition compared with others or boundaries tested & Always \\
B05 & Tasks and answers checked & Checking procedure described & Coverage, findings, and limits of check reported & Always \\
B07 & Scoring implementation checked & Scoring code inspectable & Parser validated; exploit / anti-gaming tests reported & Automated / mixed \\
B12 & Uncertainty and statistics & Sample size and score uncertainty & Power or significance reasoning reported & Always \\
B14 & Contamination assessed & Risk acknowledged; safeguard given & Contamination check / held-out procedure reported & Always \\
B18 & Replication instructions & Code, data, environment present & Independent replication reported; dependencies locked & Always \\
B28 & Execution isolated & Environment isolation described & State clearing and sandbox isolation verified & Code execution \\
\bottomrule
\end{tabularx}
\end{table}

\section{Empirical Research Design}
\label{sec:research_design}

Our empirical evaluation addresses four research questions:
\begin{itemize}[leftmargin=*]
    \item \textbf{RQ1 (Coverage):} Which review concerns do existing frameworks cover?
    \item \textbf{RQ2 (Concordance):} Do frameworks agree when applied to identical benchmarks?
    \item \textbf{RQ3 (Unstated Information):} How does each framework treat missing reporting?
    \item \textbf{RQ4 (Human-in-the-Loop Audit):} Can software engineers effectively verify coding agent reviews using the MESA audit workbench, and do they catch code defects?
\end{itemize}

\subsection{Study 1 \& 2: Model Studies Across Four Benchmarks}
In Study 1, coding agents (Claude Sonnet 5.5 through Claude Code at medium effort) applied six frameworks across four frozen benchmarks: HumanEval \cite{human_eval_repository}, TruthfulQA \cite{truthfulqa_repository}, Humanity's Last Exam (HLE) \cite{phan_etal_hle_2025}, and $\tau$-bench \cite{tau_bench_repository}. In Study 2, agents applied the MESA rubric twice to all four benchmarks. All sessions were confined to read-only tools (\texttt{list\_files}, \texttt{read\_file}) over static snapshots.

\subsection{Study 3: Human-in-the-Loop Audit Workbench}
To answer RQ4, we developed an interactive, single-page audit platform (\texttt{surveys/rubric-audit/}). Software engineering participants inspect the 37 ratings produced by Claude Sonnet 5.5 on HumanEval. The workbench presents a dual-pane interface: the left pane embeds the full repository code (\texttt{human\_eval/execution.py}, \texttt{README.md}, etc.), while the right pane presents a progressive-disclosure audit flow (Glance $\to$ Evidence Quote $\to$ Verdict: \emph{Confirmed}, \emph{Overstated}, \emph{Absence Plausible}, or \emph{Contradicted}). Participants also complete five required perception ratings on clarity, auditing ease, workload, confidence, and human-AI pairing value.

\subsection{Analysis and Metrics}
We compute exact pairwise agreement and chance-corrected agreement (Cohen's $\kappa$). For human audits, we evaluate concordance with reference checks, defect detection rates (specifically HumanEval's active execution call), and distributions across the five perception ratings.

\section{Empirical Results}
\label{sec:results}

\subsection{Coverage and Disagreements (RQ1 \& RQ2)}
No criterion is covered by all seven frameworks. Across 18 framework-benchmark sessions in Study 1, standardized verdicts agree on only 53\% of shared criteria. Across 143 pairwise disagreements, verified coding reveals that \textbf{92 of 143 discrepancies (64\%) stem from item granularity (50) or unstated-information rules (42)}, rather than conflicting interpretations of evidence. In contrast, same-framework repeats on HumanEval agreed on 19/21 (BetterBench) and 20/22 criteria (Scorecard), demonstrating that divergence reflects framework design rather than stochastic variation.

\subsection{The 10:1 Ratio of Omissions to Defects (RQ3)}
In Study 1 agent sessions, penalties for unstated information outnumbered confirmed defects by over ten to one: across 711 answers, frameworks recorded 199 answers with absent information (52 BB, 123 SC, 23 H2B, 12 ABC) and only 13 with contradicted evidence. Furthermore, 210 answers (30\%) carried the \emph{undocumented} qualifier. Five frameworks penalize omissions by design, leading to false-negative rejections of benchmarks whose only fault is concise documentation.

\subsection{Human Audit Findings and Human-Agent Concordance (RQ4)}
\noindent\textbf{Human evaluation results for Study 3:}
\pending{Insert participant cohort demographics (N=10--15 software engineers/researchers), recruitment dates, and completion denominators.}

Table~\ref{tab:human_audit_results} summarizes human auditor verification of the agent's 37 ratings on HumanEval.
\begin{table}[t]
\caption{Human auditor verification of Claude Sonnet 5.5's HumanEval review.}
\label{tab:human_audit_results}
\centering
\scriptsize
\begin{tabular}{lrrrrr}
\toprule
\textbf{Audit Stage} & \textbf{Items} & \textbf{Confirmed} & \textbf{Overstated} & \textbf{Absence Plausible} & \textbf{Contradicted} \\
\midrule
Part A (OECD Facets) & 8 & \pending{N} & \pending{N} & --- & \pending{N} \\
Part B (Rubric Questions) & 29 & \pending{N} & \pending{N} & \pending{N} & \pending{N} \\
\midrule
\textbf{Total} & \textbf{37} & \pending{Total} & \pending{Total} & \pending{Total} & \pending{Total} \\
\bottomrule
\end{tabular}
\end{table}

\noindent\textbf{Defect Detection in Code:} In Question B18, the agent rated replication materials \emph{Shown} but erroneously asserted in its reasoning that \emph{``execution is commented out by default''}. When human auditors used the workbench's code viewer to inspect \texttt{human\_eval/execution.py} line 50, \pending{insert \% of human auditors} successfully caught the active \texttt{exec(check\_program, exec\_globals)} call and flagged the agent's reasoning as \emph{Contradicted}.

\noindent\textbf{Inter-Rater Agreement:} Chance-corrected agreement between human auditors and the agent achieved \pending{insert Cohen's kappa $\kappa = X.XX$}. Human rater concordance was highest on clearly anchored criteria (B01, B06, B19) and lower on boundary distinctions (B02 user relevance vs model use).

\subsection{Auditor Perceptions and Usability Feedback}
The five perception feedback ratings demonstrated that human-agent pairing provides substantial practical benefits for requirements auditing:
\begin{itemize}[leftmargin=*]
    \item \textbf{Manageable Workload:} Average completion time was \pending{insert average minutes, e.g. 31.4 minutes}, with \pending{insert \%} rating workload as manageable or low effort.
    \item \textbf{Human-AI Pairing Utility:} \pending{insert \%} rated starting from the agent's pre-extracted evidence as \emph{noticeably faster} or \emph{substantially better} than auditing from scratch.
    \item \textbf{Review Confidence:} \pending{insert \%} expressed high confidence in the final verified benchmark review.
\end{itemize}

\section{Discussion: Implications for RE4AI Practice}
\label{sec:discussion}

\subsection{Auditing AI Benchmarks for Acceptance Decisions}
When software engineers specify that an AI model must achieve a target threshold on a benchmark to satisfy a system requirement, MESA provides the acceptance audit checklist:
\begin{enumerate}[leftmargin=*]
    \item \textbf{Construct Alignment (B01--B04):} Does the benchmark operationalize the specific ability required by the software feature?
    \item \textbf{Execution Security (B07, B28):} Does the evaluation harness execute generated code safely without active exploits?
    \item \textbf{Uncertainty Bounds (B12):} Are scores backed by confidence intervals or power analyses to ensure differences between models are statistically discernible?
\end{enumerate}

\subsection{Requirements Traceability Matrix for AI Components}
In classical requirements engineering, a Requirements Traceability Matrix (RTM) links high-level business goals to software specifications and verification test cases. For AI-enabled systems, MESA extends this traceability chain:
\[
\text{System Requirement} \longrightarrow \text{Construct Definition (B01)} \longrightarrow \text{Benchmark Tasks (B03--B04)} \longrightarrow \text{Audit Evidence (MESA)}
\]
By establishing explicit traceability between system acceptance criteria and benchmark audit evidence, engineering teams can verify whether third-party model claims provide sound acceptance verification.

\subsection{Human-Agent Teaming in Requirements Verification}
Our findings show that neither agents nor humans alone provide an ideal audit: coding agents rapidly scan text and extract locators, but miss active execution calls and overstate design facets; humans provide critical verification judgment but face severe fatigue if forced to read 35-page papers from scratch. Pairing human engineers with agent-extracted evidence cards resolves this dilemma, cutting review time to $\sim$30 minutes while maintaining rigorous oversight.

\subsection{Decoupling Omissions from Defects in AI Procurement}
In safety-critical procurement (e.g. EU AI Act conformity \cite{eu_ai_act_2024}), conflating unstated documentation with empirical failure leads to erroneous supplier rejections. MESA's \emph{Undocumented} qualifier ensures that missing documentation triggers targeted inquiries to vendors rather than false-positive defect reports.

\section{Threats to Validity}
\label{sec:threats}

\subsection{Construct Validity}
The 30 criteria represent interpretive syntheses of 280 framework items. While all mappings and reconciliation notes are published, alternative groupings are possible. Standardization rules were registered before sessions.

\subsection{Internal Validity}
The MESA rubric was derived before study sessions and excluded from Study 1 agreement counts. In Study 3, participants reviewed agent-extracted citations rather than discovering citations from scratch, which may anchor auditor judgments to the agent's search paths.

\subsection{External Validity}
Our findings are bounded by four benchmarks and one coding agent model family (Claude Sonnet 5.5). In Study 3, HumanEval represents code synthesis; transferring these findings to vision or multimodal agent benchmarks requires future evaluation.

\subsection{Conclusion Validity and Reliability}
Automated validation scripts ensure all counts and tables are deterministic. For human evaluations, we report exact denominators, chance-corrected agreement ($\kappa$), and qualitative comment coding.

\section{Conclusion and Future Work}
\label{sec:conclusion}

MESA bridges the divide among existing benchmark review frameworks by establishing a shared vocabulary of 30 criteria, standardized verdicts that isolate missing documentation from verified defects, and an evidence-grounded rubric. Our empirical studies show that cross-framework disagreements stem overwhelmingly from scoring rules rather than evidence differences, and that negative scores penalize unstated practices ten times more often than confirmed defects. Through our human-in-the-loop audit platform, software engineers verified agent reviews in $\sim$32 minutes while catching code execution flaws that static inspection missed. Future work will expand the human audit platform across multimodal benchmarks and develop automated repair heuristics for benchmark harnesses.

\section*{Data Availability}
Study materials, benchmark snapshots, the interactive audit application, and analysis scripts are available at an anonymous repository: \url{https://anonymous.4open.science/r/mesa-audit-placeholder}.

\bibliographystyle{splncs04}
\bibliography{references}

\end{document}
